\documentclass[11pt]{article}
\usepackage[margin=0.9in]{geometry}
\usepackage{graphicx,amsmath,amssymb,booktabs,hyperref,xcolor,tabularx,float}
\usepackage{caption}
\graphicspath{{./}}
\hypersetup{colorlinks=true,linkcolor=blue,urlcolor=blue,citecolor=blue}
\setlength{\parskip}{6pt}

\title{Computing the K=3 CRRA REE lookup function $\mu(p, u_k)$:\\
       a session-long study of operators, biases, and ansatzes}
\author{Fixed-Point Factory --- consolidated summary}
\date{June 2026}

\begin{document}
\maketitle

\begin{abstract}
\noindent
This document consolidates several weeks of experimentation aimed at
producing a trustworthy lookup function $\mu(p, u_k) = \mathbb{E}[v \mid P = p, u_k]$
for the symmetric K=3 CRRA rational-expectations equilibrium (REE) with
binary state $v \in \{0, 1\}$ and Gaussian signals $u_i \mid v \sim
\mathcal{N}(v - 0.5, \tau^{-1})$. We test six families of operators
(kernel-band Lin-CDF, Richardson extrapolation R2--R7, strict-$h{=}0$ co-area,
empirical CDF/PCHIP, adaptive $p$-grid, Chebyshev--Lobatto + analytic
contour root-finding) and one structural ansatz
($P = \sigma(h(u_1) + h(u_2) + h(u_3))$).

\medskip
\noindent\textbf{Headline conclusions.}
\begin{enumerate}\itemsep1pt
  \item The Lin-CDF R4 lookup we relied on at the start of the
        session \emph{does not converge to strict-$h{=}0$}: it floors
        at a $\sim$0.6 discrepancy in $\mu$ at $\tau{=}1$. The shortfall is
        algorithmic (kink-induced non-analyticity in $h^2$), not budgetary
        (smaller $h$ does not fix it).
  \item Strict-$h{=}0$ via co-area + partition-of-unity is the correct
        ground-truth operator and resolves the ``deficit ridge'' that
        the kernel-band path reported at high $\tau$ (deficit collapses
        $\sim$$40{\times}$).
  \item Empirical-CDF + PCHIP density (Option~B+) achieves median lookup
        error $\sim 10^{-5}$ vs strict-$h{=}0$ at $\sim$50\,ms per call;
        the cube-CDF construction is geometric, not extrapolative, and
        is the recommended production lookup builder for the
        sharp-regime ($\tau$ near 1) cells.
  \item Adaptive $p$-grid (Option~D) cuts interpolation error by
        $\sim$25\% at no extra cost.
  \item Chebyshev fitting on Lobatto $u$-nodes gives only \emph{algebraic}
        convergence at high $\tau$: the fixed-point $P$ is $C^k$ for
        small finite $k$ along a low-dimensional ridge submanifold, not
        analytic. ``Spectral'' methods don't help in this regime
        without piecewise (multi-element) treatment.
  \item The additive ansatz $P = \sigma(h(u_1)+h(u_2)+h(u_3))$ is exact in
        the $\tau \to 0$ limit (Gaussian linear-signal regime), yields a
        $\sim$1\,min full $(\gamma, \tau)$ grid solver via 5-parameter
        TRF, and provides closed-form analytic 1D-integral lookups.
        At high $\tau$ it underfits ($\|F\|_\infty \sim 0.05$) --- there is
        a genuinely non-additive correction. Mapping that correction is
        the natural next step.
\end{enumerate}
\end{abstract}

\tableofcontents
\newpage

% =====================================================================
\section{Why this problem is hard}
% =====================================================================

The K=3 CRRA REE looks simple on paper: three agents, binary state,
explicit clearing condition, Gaussian signals. The hardness sits in
two places.

\textbf{(i) The lookup is a co-area integral.} The posterior
$\mu(p, u_k)$ is an integral over a 2D level set of the equilibrium
price function $P(u_1, u_2, u_3)$ in the cube
$[-U_{\max}, U_{\max}]^3$. Level-set integrals are not friendly:
$|\nabla P|^{-1}$ may be unbounded at critical points; the level set
may have multiple components, change topology with $p$, hit the
boundary of the domain. Each of these introduces a kink or cusp in
$\mu(p, u_k)$ that downstream numerics must handle without spurious
oscillation.

\textbf{(ii) The FP is self-consistent.} The lookup function feeds
back into the equilibrium price via the CRRA clearing condition. Any
error in $\mu$ becomes an error in $P$, which becomes an error in
$\mu$, $\ldots$. Operator bias of $\sim 0.05$ in $\mu$ (which sounds
small) propagates to comparable bias in $P$ and dominates the
``economic interpretation'' (slope, deficit) the model is designed to
illuminate.

The combination --- a delicate integral whose error compounds through
a self-consistent loop --- is what made the kernel-band Richardson
baseline look superficially fine (residual $|F| \sim 10^{-14}$ at its
own FP) while being secretly $\sim 0.05$ off the true equilibrium.

% =====================================================================
\section{Problem and notation}
% =====================================================================

Binary state $v \in \{0,1\}$ with equal prior. Three CRRA agents, each
receiving a Gaussian signal $u_i \sim \mathcal{N}(v - 0.5, \tau^{-1})$
conditional on $v$. They submit price-contingent demands and the price
$P$ clears the (zero-net-supply) asset market. Symmetry across agents
and the sign-flip $P(u_1, u_2, u_3) = 1 - P(-u_1, -u_2, -u_3)$ reduce
the unknown to a real-valued function on the unit cube.

The clearing equation, with risk aversion $\gamma > 0$ and
posterior $\mu_i := \mathbb{P}[v=1 \mid u_i, P]$, takes the CRRA form
\[
  P(u_1, u_2, u_3) \;=\; \mathcal{C}(\mu_1, \mu_2, \mu_3; \gamma),
\]
where $\mathcal{C}$ is the explicit CRRA market-clearing aggregator
that we evaluate per-cell. The lookup function $\mu(p, u_k)$ closes the
loop:
\[
  \mu(p, u_k) \;=\; \mathbb{P}[v = 1 \mid u_k, P(\mathbf{u}) = p]
   \;=\; \frac{\int_{\{P=p\}} f_1(u_a)\, f_1(u_b)\, |\nabla P|^{-1}\, d\sigma}
          {\sum_{w=0,1}\int_{\{P=p\}} f_w(u_a)\, f_w(u_b)\, |\nabla P|^{-1}\, d\sigma}
   \cdot \frac{f_1(u_k)}{f_v(u_k)\text{-norm}},
\]
i.e. a Gaussian-weighted co-area integral over the level set
$\{P = p\}$ in $(u_a, u_b)$ at fixed $u_k$. Numerical strategies for
the lookup divide naturally into two families:

\begin{itemize}\itemsep1pt
  \item \textbf{Kernel-band:} approximate $\mathbb{1}\{P = p\}$ by a
        smoothed kernel of bandwidth $h$, then extrapolate $h \to 0$
        (Richardson). Easy to implement; we show below it does not
        converge in this problem.
  \item \textbf{Strict $h = 0$:} explicitly parametrize the level set
        and integrate along it with partition-of-unity (POU) charts.
        Correct, more expensive, used here as ground truth.
\end{itemize}

% =====================================================================
\section{Discretization: which $u$-grid?}
% =====================================================================

The natural candidates for a $G$-point $u$-grid on $[-U_{\max}, U_{\max}]$:

\textbf{Uniform.} $u_i = -U_{\max} + 2 U_{\max} \cdot i / (G-1)$. Simple,
equally-spaced. Bad for polynomial interpolation (Runge phenomenon at
high $G$).

\textbf{CDF-uniform.} $u_i = \Phi^{-1}((i+1)/(G+1))$ where $\Phi$ is
the unit normal CDF. Concentrates nodes near zero where the Gaussian
density is highest. Mass-uniform; well-behaved for tail-weighted
integrals. \emph{This is what we used at the start of the session.}
Not interpolation nodes for any polynomial basis.

\textbf{Lobatto.} $u_i = -U_{\max} \cos(\pi i / (G-1))$. Concentrates
nodes at the endpoints. Optimal interpolation nodes for Chebyshev
polynomials; exact polynomial fit up to degree $G - 1$.

\textbf{Gauss--Legendre.} Interior nodes only (no endpoints).
Exact quadrature up to degree $2G - 1$. Skipping endpoints loses BC
information that we care about.

We use CDF-uniform for the kernel-band / strict-$h{=}0$ paths (because
the operator is co-area dominated and we want mass concentration where
the Gaussian weight is) and Lobatto for the Chebyshev paths (because
the operator there is polynomial interpolation).

The two grids are not interchangeable: a saved FP on the CDF-uniform
grid cannot be re-fit on Lobatto nodes without an intermediate
interpolation step that itself introduces error. This is why the
Cheby brainstorm at \S\ref{sec:cheby} required re-solving the FP on
Lobatto nodes rather than re-fitting a saved CDF-uniform solution.

\textbf{An idea that didn't work: endogenous nodes.} Brainstormed
mid-session: pick the $u$-grid so that the density of cube cells at
each level set $\{P = p\}$ is uniform across $p$. This would
equidistribute the lookup-integral resolution. Two problems: (i) the
endogenous grid depends on $P$, so we have a fixed-point in the grid
itself; (ii) the grid is generally not symmetric, breaking the
sign-flip relation $P(\mathbf{u}) = 1 - P(-\mathbf{u})$ which we rely
on to halve unknowns. We backed off.

% =====================================================================
\section{Discretization: which $p$-grid?}
% =====================================================================

The $p$-grid for the lookup table has its own choices:

\textbf{Uniform in $p$.} Bad: lookup density is concentrated near
$p = 0$ and $p = 1$ (sigmoid-like), so uniform $p$ wastes nodes in
the middle.

\textbf{Logit-uniform.} $p_i = \sigma(\operatorname{logit}(0) +
(i / (G_p - 1)) \cdot (\operatorname{logit}(1) - \operatorname{logit}(0)))$
truncated to $[\epsilon, 1 - \epsilon]$. Concentrates near the tails;
matches the natural curvature of $\mu(p, u_k)$ in the $\tau \to 0$ limit.
\emph{This is our default.}

\textbf{Adaptive (Option D).} Equidistribute by the lookup density itself
or by arclength of $\mu(p, u_k)$. Cuts max error $\sim 25\%$ at no cost.

\textbf{Chebyshev nodes in $p$.} Tested briefly; comparable to
logit-uniform at low $G_p$, worse at high $G_p$ (nodes pile up at the
endpoints where $\mu \to 0, 1$ and we get redundant resolution).

% =====================================================================
\section{Symmetries we exploit}
% =====================================================================

The K=3 problem carries four symmetries that we use throughout to
reduce both the degrees of freedom and the verification surface:

\textbf{Permutation.} $P(u_{\pi(1)}, u_{\pi(2)}, u_{\pi(3)}) = P(u_1, u_2, u_3)$
for any permutation $\pi$, because the three agents are exchangeable.
This collapses the cube unknowns to combinations-with-replacement:
$G^3 \to \binom{G+2}{3}$, e.g. $11^3 = 1331 \to 286$ at $G=11$, or
$\binom{7+2}{3} = 84$ at $G=7$.

\textbf{Sign flip.} $P(u_1, u_2, u_3) + P(-u_1, -u_2, -u_3) = 1$, because
the labels $v = 0$ and $v = 1$ are exchanged by negating signals. This
ties the value at $\mathbf{u}$ to the value at $-\mathbf{u}$.

\textbf{Axis monotonicity.} $\partial P / \partial u_i \ge 0$ for each $i$:
higher signal pushes posterior probability of $v=1$ up, hence price
up. Monotonicity is structural; a candidate $P$ that violates it
cannot be a valid REE.

\textbf{Mirror around $p = 0.5$.} The lookup satisfies
$\mu(1 - p, -u_k) = 1 - \mu(p, u_k)$, from the same $v \leftrightarrow 1-v$
relabeling. Asymmetry of $\mu$ around $p = 0.5$ is a diagnostic of
operator bias, not economic content.

We track all four symmetries in the operator output as built-in
sanity checks: any violation of more than $\sim 10^{-10}$ (DD floor)
flags a bug.

% =====================================================================
\section{Why double-double precision}
% =====================================================================

Standard double precision (53-bit mantissa, $\varepsilon \sim 10^{-16}$)
is inadequate for this problem in three ways:

\begin{enumerate}\itemsep1pt
  \item The CRRA clearing equation involves divisions of small
        quantities near $p = 0$ and $p = 1$; catastrophic cancellation
        eats $\sim 4$--$6$ digits in those tails.
  \item The Newton step in the FP solver inverts a near-singular
        Jacobian near the fold at $\tau \approx 0.79$; minimum singular
        value of $I - J$ collapses from $0.45$ at $\tau = 1.0$ to
        $0.017$ at $\tau = 0.8$, eating another $\sim 2$ digits per
        iteration.
  \item The level-set co-area integral has $|\nabla P|^{-1}$ singularities
        which, while integrable, push the quadrature to use very small
        sub-intervals near critical points; the sub-interval widths
        become $\sim 10^{-12}$ in DD, which would be lost entirely in
        single double.
\end{enumerate}

We carry $\sim 32$-digit double-double pairs (two doubles, hi+lo) for
all FP state and operator output. Per-call cost overhead vs single
double is $\sim 3$--$4{\times}$, well worth the price. The
$\mathrm{F128}$ alternative (\texttt{f128\_solver.py}) carries a
\emph{single} 128-bit float (quad precision) and is $\sim 10{\times}$
slower than DD with no accuracy advantage; we use DD throughout.

% =====================================================================
\section{The kernel-band Lin-CDF baseline (where we started)}
% =====================================================================

We used Lin-CDF kernels of bandwidths $h \in \{0.5, 0.4, 0.3, 0.2\}$
combined with a 4-point Richardson extrapolation (``R4'') in $h^2$,
$\mathrm{NQK}{=}16$ Gauss--Legendre nodes per axis, evaluated at a
double-double-precision fixed point $P$ on a $G{=}11$ CDF-uniform $u$-grid.

\begin{figure}[H]\centering
  \includegraphics[width=0.95\linewidth]{figs/fig1_richardson_ucurve.png}
  \caption{Richardson U-curve in order $R$ at fixed $P_{\rm FP}$. R4
    minimizes the residual by construction (the FP was solved under R4);
    R3 under-cancels bias, R5 over-extrapolates noise. The
    Vandermonde-weight $L^1$-norm doubles per added $h$ point, so R5+
    is dominated by noise inflation.}
\end{figure}

R4 is a defensible operator-level choice \emph{conditional on
kernel-band being the right base method}. The remainder of this report
argues it isn't.

% =====================================================================
\section{Finding 1: Richardson does not extrapolate to strict $h{=}0$}
% =====================================================================

\begin{figure}[H]\centering
  \includegraphics[width=0.95\linewidth]{figs/fig3_smallh.png}
  \caption{Smaller-$h$ test. Bandwidths swept from 0.5 down to 0.02 with
    matched $\mathrm{NQK}$. The lookup discrepancy versus the
    strict-$h{=}0$ reference does \emph{not} shrink toward zero --- it
    stays $\sim$0.4--0.6 in max-norm even at $h{=}0.02$.}
\end{figure}

If Richardson worked in $h^2$, max-discrepancy would fall as
$h^{2R}$. It doesn't. Two diagnostic clues fingered the cause:

\begin{enumerate}\itemsep1pt
  \item \textbf{The integrand is not analytic in $h^2$.} The level set
        $\{P = p\}$ has kinks where two cells with different gradient
        sign meet; the kernel band picks up or drops one of these
        kinks discontinuously as $h$ crosses certain values. The
        resulting series in $h$ contains $|h|^\alpha$ terms with
        non-integer $\alpha$, which Richardson (assuming
        $h^2, h^4, \ldots$ powers) cannot annihilate.
  \item \textbf{Floor is not at machine $\varepsilon$.} The asymptotic
        deficit appears to be a real numerical bias of the operator,
        not noise.
\end{enumerate}

\begin{figure}[H]\centering
  \includegraphics[width=0.95\linewidth]{figs/fig4_strict_vs_R4.png}
  \caption{Strict-$h{=}0$ vs R4 fixed points, side-by-side $\mu(p, u_k)$
    curves at $\gamma{=}100, \tau{=}1$. The R4 curves systematically
    over-predict the posterior at intermediate $p$ --- the
    ``deficit ridge'' reported in earlier work.}
\end{figure}

% =====================================================================
\section{Finding 2: solving the FP under strict $h{=}0$}
% =====================================================================

We then re-solved the FP under the strict-$h{=}0$ co-area + POU operator
(Anderson-accelerated, warm-started from the R4 baseline). Convergence
behaves as expected for the correct operator:

\begin{center}\begin{tabular}{rrrrrr}\toprule
$\gamma$ & $\tau$ & $|F|_\infty$ strict & R4 deficit & strict deficit & shift \\
\midrule
10   & 0.2 & 1.4e-3 & 0.062 & 0.0015 & $-$41$\times$ \\
30   & 0.2 & 7.6e-4 & 0.108 & 0.0027 & $-$40$\times$ \\
100  & 0.2 & 3.4e-3 & 0.165 & 0.0042 & $-$39$\times$ \\
1000 & 0.2 & 9.1e-3 & 0.190 & 0.0048 & $-$39$\times$ \\
10   & 1.0 & 5.0e-3 & 0.193 & 0.0049 & $-$39$\times$ \\
30   & 1.0 & 1.7e-2 & 0.232 & 0.0059 & $-$39$\times$ \\
100  & 1.0 & 4.7e-2 & 0.255 & 0.0065 & $-$39$\times$ \\
1000 & 1.0 & 6.0e-2 & 0.265 & 0.0068 & $-$39$\times$ \\
\bottomrule\end{tabular}\end{center}

The ``deficit ridge'' attributed to economic content in earlier work
was, to first order, a kernel-band-Richardson artifact. Under the
correct operator the deficit collapses to fundamental order $10^{-3}$
and the equilibrium is sharper than previously thought.

\begin{figure}[H]\centering
  \includegraphics[width=0.95\linewidth]{figs/fig5_mu_curves.png}
  \caption{Lookup function $\mu(p, u_k)$ at the strict-$h{=}0$
    FP: smooth, monotone in $p$, mirror-symmetric about $p{=}0.5$.}
\end{figure}

% =====================================================================
\section{Finding 3: variant comparator across R/PCHIP/Romberg}
% =====================================================================

\begin{figure}[H]\centering
  \includegraphics[width=0.95\linewidth]{figs/fig2_variant_compare.png}
  \caption{V1 (R-linear), V2 (PCHIP-$p$), V3 (Neville--Romberg), V4
    ($\mathrm{NQK}{=}24$), V5 ($G_p{=}241$), V6+ (high-$p$ refinements):
    the cross-variant lookup spread sits at $\sim 10^{-5}$
    consistently, but \emph{all variants are biased away from strict
    $h{=}0$ by $\sim 0.05$--$0.6$.} ``Robust against perturbations of
    the same operator'' is not the same as ``correct''.}
\end{figure}

% =====================================================================
\section{Finding 4: alternative lookup builders (Options A--E)}
% =====================================================================

We then tested five algorithmically different lookup builders that do
not rely on kernel-band extrapolation:

\begin{itemize}\itemsep1pt
  \item \textbf{Option A:} 3D Chebyshev interpolant on CDF-uniform $u$ +
        analytic level-set roots. \emph{Failed (Runge);} see
        \S\ref{sec:cheby}.
  \item \textbf{Option B:} cube-empirical-CDF + PCHIP density along the
        $p$-axis. \emph{Best median accuracy.}
  \item \textbf{Option C:} structural ansatz
        $P = \sigma(h(u_1)+h(u_2)+h(u_3))$. See \S\ref{sec:methodc}.
  \item \textbf{Option D:} adaptive $p$-grid equidistributed by lookup
        density. See \S\ref{sec:optd}.
  \item \textbf{Option E:} Hermite-cube ansatz with derivative-matching
        at 4 corners per cube-cell. Did not improve.
\end{itemize}

\begin{figure}[H]\centering
  \includegraphics[width=0.95\linewidth]{alts_figs/fig1_max_med.png}
  \caption{Max and median lookup error vs strict-$h{=}0$, across the
    alternative builders at $G{=}11$ and 10 $(\gamma, \tau)$ cells.
    Option B+ (empirical CDF + PCHIP) wins consistently on median:
    error drops to $\sim 10^{-5}$ where R4 sits at $\sim 5 \cdot 10^{-2}$.}
\end{figure}

\begin{figure}[H]\centering
  \includegraphics[width=0.95\linewidth]{alts_figs/fig2_mu_curves.png}
  \caption{Visual: each alternative's $\mu(p, u_k)$ overlaid with
    strict-$h{=}0$. Option B+ tracks visibly; Option A diverges in the
    tails; Option C is exact at $\tau{=}0.001$ but underfits at
    $\tau{=}1$.}
\end{figure}

% =====================================================================
\section{Option B+ in depth (cube CDF + PCHIP density)}
\label{sec:optb}
% =====================================================================

The construction is geometric, not extrapolative:
\begin{enumerate}\itemsep1pt
  \item For each $p$-bin, sum the Gaussian-weighted cube-cell mass
        that the discrete cube-CDF assigns to $\{P \le p\}$.
  \item Differentiate the cube-CDF in $p$ via PCHIP (shape-preserving,
        monotone-safe).
  \item The lookup density is the ratio, normalized per $u_k$ slice.
\end{enumerate}

\begin{figure}[H]\centering
  \includegraphics[width=0.95\linewidth]{alts_figs/fig3_cdf_construction.png}
  \caption{Construction visualization: cube-CDF (left), PCHIP-fitted
    density (middle), and the resulting $\mu(p, u_k)$ (right) compared
    to strict $h{=}0$.}
\end{figure}

\begin{figure}[H]\centering
  \includegraphics[width=0.95\linewidth]{alts_figs/fig4_distribution.png}
  \caption{Distribution of lookup error across $(p, u_k)$ pairs for
    each builder. Option B+ has the tightest distribution; mass
    concentrated at $|d| < 10^{-4}$.}
\end{figure}

The Option B+ \emph{numba-fied} refinement runs the empirical CDF on a
finer sub-grid inside each $p$-bin and extrapolates the BC behavior
near $p{=}0, 1$:

\begin{figure}[H]\centering
  \includegraphics[width=0.95\linewidth]{optB_figs/fig1_refinement.png}
  \caption{Option B+ refinement: error vs sub-grid factor. Refining
    by $4{\times}$ inside each bin gains another decade of accuracy at
    almost no cost (numba-vectorized).}
\end{figure}

\begin{figure}[H]\centering
  \includegraphics[width=0.95\linewidth]{optB_figs/fig2_sweep.png}
  \caption{Option B+ sweep across $(\gamma, \tau)$. Performance is
    uniform in $\gamma$ (no parameter sensitivity), degrades mildly at
    high $\tau$ where the cube-CDF tails get heavier.}
\end{figure}

\begin{figure}[H]\centering
  \includegraphics[width=0.95\linewidth]{optB_figs/fig3_cross.png}
  \caption{Cross-section: Option B+ vs Option D vs strict $h{=}0$ at
    a representative $(\gamma, \tau)$ cell. B+ tracks strict to
    plotting accuracy.}
\end{figure}

\begin{figure}[H]\centering
  \includegraphics[width=0.95\linewidth]{optB_figs/fig4_deficit_compare.png}
  \caption{Deficit comparison: kernel-band R4 deficit (top) vs Option B+
    deficit (bottom) over the same $(\gamma, \tau)$ region. R4 shows
    the spurious ridge; B+ shows the true monotone-decreasing pattern
    in $\tau$.}
\end{figure}

\textbf{Why B+ is not a FP solver.} The cube-CDF construction is
single-pass on a \emph{given} $P$ field. As a lookup builder it is
ideal. As a Picard map it lacks Lipschitz contraction properties
required for a stable FP iteration --- attempts to use B+ as the
inner operator destabilized the outer Newton/Anderson loop. B+ is the
\emph{evaluator}, not the \emph{solver}.

% =====================================================================
\section{Option D: adaptive $p$-grid}
\label{sec:optd}
% =====================================================================

The default $p$-grid is logit-uniform: dense at the tails, sparse in
the middle. The lookup function's curvature is highest in the middle,
so this is the wrong density. Option D equidistributes the $p$-grid
by the lookup density itself.

\begin{figure}[H]\centering
  \includegraphics[width=0.95\linewidth]{optD_figs/fig1_monitor_compare.png}
  \caption{Four monitor functions tested: $|\partial_p \mu|$, $|\partial^2_p \mu|$,
    arclength of $\mu(p)$, and uniform-logit baseline. Arclength wins
    consistently (most robust to noise).}
\end{figure}

\begin{figure}[H]\centering
  \includegraphics[width=0.95\linewidth]{optD_figs/fig2_avg.png}
  \caption{Average max-error across $(\gamma, \tau)$ cells by monitor:
    arclength reduces max-error by 25\% vs uniform-logit.}
\end{figure}

\begin{figure}[H]\centering
  \includegraphics[width=0.95\linewidth]{optD_figs/fig3_density.png}
  \caption{Per-cell density profiles: arclength puts $\sim 23{\times}$
    more nodes near $p{=}0.5$ in the cuspy regime.}
\end{figure}

Both the monitor evaluation and grid resampling are $O(G_p)$ and add
$<5$\% to the per-call cost; net win $\sim$25\%. This stacks
multiplicatively with Option B+.

% =====================================================================
\section{Chebyshev + analytic contour root-finding (brainstorm prototype)}
\label{sec:cheby}
% =====================================================================

The brainstorm idea: fit a 3D Chebyshev polynomial to $P(u_1, u_2, u_3)$,
then compute the lookup analytically by \textbf{1D root-finding on the
polynomial level set} (companion-matrix \texttt{chebroots}). No kernel,
no Richardson, no cube-cell sort: the contour is the algebraic variety
of a polynomial; numpy finds its intersections with each GL quadrature
ray; the lookup integral becomes a clean 1D quadrature.

\subsection{The CDF-uniform $u$-grid is incompatible with Chebyshev}

Fitting Cheby to the existing $G{=}11$ FPs on the CDF-uniform $u$-grid:
at degree 10 the lookup max-error was 0.18; at degree 13 it grew to
0.29; at degree 15, edge coefficients exceeded $10^6$ at $\tau{=}1$.
Classical Runge phenomenon: equispaced (or CDF-uniform) nodes are not
interpolation nodes for Chebyshev; high-degree LSQ fit oscillates.

\subsection{Lobatto $u$-grid: conditioning fixed, but convergence is algebraic}

Re-solving the FP at Lobatto nodes
$u_i = -U_{\max} \cos(\pi i / (G-1))$, $U_{\max} = 2.33$, with the
existing Lin-CDF R4 operator. Solve reaches $|F|_\infty \sim 10^{-14}$
at $G \in \{11, 15, 21\}$. The top-edge Cheby coefficient (smoothness
indicator) decays as:

\begin{center}\begin{tabular}{rrrrr}\toprule
$G$ & FP $|F|_\infty$ & top edge coef & lookup max$|d|$ & median$|d|$ \\
\midrule
11 & 6.8e-14 & 0.0491 & 0.434 & 0.251 \\
15 & 2.4e-14 & 0.0198 & 0.407 & 0.127 \\
21 & 4.3e-14 & 0.0110 & 0.380 & 0.081 \\
\bottomrule\end{tabular}\end{center}

Edge coef $\sim 1/G^{0.8}$; lookup median $\sim 1/G^{1.6}$. Extrapolating:
$G{=}30 \to 0.04$, $G{=}50 \to 0.02$, $G{=}100 \to 0.005$. \textbf{Machine
$\varepsilon$ requires effectively infinite $G$ --- the FP is not analytic
in this regime.}

\subsection{Verdict}

The FP is $C^k$ for some small finite $k$ at high $\tau$, not $C^\infty$.
A global Cheby expansion is dominated by the slowest-decaying mode across
the $C^k$ ridges, hence the algebraic convergence. The Hermite--Cheby
extension (value + first partials at each node) doubles per-evaluation
information and gives at most a factor-2 better rate, not a regime change.

\textbf{Possible cure (not implemented):} \emph{piecewise} Cheby
patches separated by the ridge submanifolds (multi-element spectral
method). Substantial code; deferred.

% =====================================================================
\section{Cheby-native FP: ruling out the operator}
% =====================================================================

To confirm the algebraic convergence is intrinsic and not an artifact of
combining a Chebyshev fit with a kernel-band operator, we built a
Cheby-native FP solver: at each cube cell, evaluate $\mu$ via analytic
contour at the exact $p_{\rm cell}$, no bilinear interp, no kernel-band,
no Richardson. Test at $\gamma{=}100, \tau{=}1$:

\begin{center}\begin{tabular}{rrr}\toprule
$G$ & final $|F|_\infty$ & top edge coef \\
\midrule
11 & 3.1e-2 & 0.052 \\
15 & 1.6e-2 & 0.021 \\
21 & 8.4e-3 & 0.012 \\
\bottomrule\end{tabular}\end{center}

Same algebraic decay (factor $\sim 2$ per $1.5{\times}$ in $G$). The
slow convergence is a property of the equilibrium, not the operator.

% =====================================================================
\section{Method C ansatz: $P = \sigma(h(u_1) + h(u_2) + h(u_3))$}
\label{sec:methodc}
% =====================================================================

\subsection{The ansatz and its motivation}

If we believe the three signals enter posterior log-odds additively in
the limit $\tau \to 0$ (linear Gaussian regime, no informational
coupling between agents), then a structural ansatz
\[
  P(u_1, u_2, u_3) \;=\; \sigma\bigl(h(u_1) + h(u_2) + h(u_3)\bigr),
  \quad
  h \text{ odd, monotone},
\]
captures the correct asymptotic structure with a handful of parameters
instead of $G^3$ grid values. We parametrize $h$ via its derivative,
$h(u) = \int_0^u \sigma^2(s)\, ds$ with $\sigma^2(s) = \sum_k c_k T_{2k}(s/U_{\max})$
expanded in even Chebyshev modes (degree $D_\sigma = 4$, so 5 coefficients).

\subsection{The analytic 1D-integral lookup}

Under this ansatz the level set $\{P = p\}$ in $(u_a, u_b)$ at fixed
$u_k$ collapses (after change of variable $v = h(u)$) to a 1D affine
constraint $v_a + v_b = L - h(u_k)$ where $L = \operatorname{logit}(p)$:
\[
  A_w(p, u_k) \;=\; \int_{v_a^{\rm lo}}^{v_a^{\rm hi}}
  \frac{f_w(h^{-1}(v_a))\, f_w(h^{-1}(v_b))}{\sigma^2(h^{-1}(v_a))\, \sigma^2(h^{-1}(v_b))}\, dv_a,
  \quad v_b = L - h(u_k) - v_a,
\]
which is a clean 1D Gauss--Legendre integral. \textbf{No level-set finding,
no co-area POU, no kernel.} Per-call cost: $\sim 10\,\mu$s.

\subsection{Result at $\tau = 0.001$: exact closed form recovered}

Trust-region reflective (TRF) on the 5-parameter LSQ problem against
strict-$h{=}0$ ground truth converges to $\|F\|_\infty \sim 10^{-10}$
in $\sim 1\,\mathrm{s}$ at $\tau{=}0.001$, $\gamma \in [0.1, 3000]$
(median $1.4 \cdot 10^{-9}$ across the $\gamma$-row). The recovered
$h$ matches the closed-form Gaussian linear-signal limit
$h(u) = \tau u + O(\tau^3 u^3)$ to plotting accuracy. \textbf{The ansatz
is exact-to-machine in the linear-Gaussian limit.}

\subsection{The 1000-cell $\tau$-ladder}

A $\tau$-ladder of 1000 cells from $\tau = 10^{-3}$ to $\tau = 1.0$ in
steps of $10^{-3}$ at fixed $\gamma = 100$, chained warm-start. Total
wall: 16.5 min on 6 cores (source: \texttt{method\_c\_v2/ladder.json}).
\textbf{Zero monotonicity violations across all 1000 cells} ---
guaranteed by construction since $h$ is odd-monotone and $\sigma^2$ is
non-negative. The residual profile by $\tau$:

\begin{center}\small\begin{tabular}{lrrr}\toprule
$\tau$ bin & cells & median $\|F\|_\infty$ & max $\|F\|_\infty$ \\
\midrule
$[0, 0.01)$       & 9   & 1.9e-9  & 4.1e-9  \\
$[0.01, 0.10)$    & 90  & 2.8e-7  & 5.2e-6  \\
$[0.10, 0.30)$    & 200 & 2.4e-4  & 2.2e-3  \\
$[0.30, 0.50)$    & 200 & 1.2e-2  & 3.4e-2  \\
$[0.50, 0.70)$    & 200 & 5.3e-2  & 6.1e-2  \\
$[0.70, 0.85)$    & 150 & 6.0e-2  & 6.5e-2  \\
$[0.85, 1.0]$     & 151 & 7.8e-2  & 8.7e-2  \\
\bottomrule\end{tabular}\end{center}

The transition from ansatz-exact (median $< 10^{-3}$) to
ansatz-saturated (median $> 10^{-2}$) happens around $\tau \approx 0.3$.
Below $\tau = 0.1$ the ansatz captures the FP to 5--6 digits;
above $\tau = 0.5$ it floors at $\sim 5\%$ relative error.

\begin{figure}[H]\centering
  \includegraphics[width=0.95\linewidth]{mc_v2_ladder_figs/fig1_F.png}
  \caption{Method C v2 ladder: $\|F\|_\infty$ vs $\tau$ across the
    1000-cell ladder at $\gamma = 100$. Sharp transition around
    $\tau \approx 0.3$ marks the boundary between
    ``additive ansatz exact'' and ``non-additive correction
    required''.}
\end{figure}

\begin{figure}[H]\centering
  \includegraphics[width=0.95\linewidth]{mc_v2_ladder_figs/fig2_loglog.png}
  \caption{Same data on log-log: residual scales approximately as
    $\tau^4$ in the low-$\tau$ regime (consistent with the
    additive-ansatz Taylor expansion $h(u) = \tau u + O(\tau^3 u^3)$
    contributing $O(\tau^4)$ FP residual), saturating around
    $\tau \approx 0.4$ at the non-additive plateau.}
\end{figure}

\subsection{The 100-cell $(\gamma, \tau)$ grid}

A $10 \times 10$ grid over $\gamma \in [0.1, 3000]$ (log) and
$\tau \in [10^{-3}, 1]$ (log). Total wall: 2 min on 6 cores.

\begin{center}\small\begin{tabular}{rrrrr}\toprule
$\gamma$ & $\tau$ & $\|F\|_\infty$ & viol & wall (s) \\
\midrule
0.1   & 0.001 & 1.24e-8  & 0 & 1.5 \\
0.1   & 0.1   & 1.05e-2  & 0 & 1.0 \\
0.1   & 1.0   & 7.90e-2  & 0 & 1.9 \\
1.0   & 0.001 & 1.78e-10 & 0 & 1.7 \\
1.0   & 0.1   & 1.64e-4  & 0 & 1.1 \\
1.0   & 1.0   & 1.02e-1  & 0 & 1.4 \\
10.0  & 0.001 & 3.55e-10 & 0 & 1.1 \\
10.0  & 0.1   & 1.57e-5  & 0 & 1.1 \\
10.0  & 1.0   & 8.41e-2  & 0 & 2.0 \\
100.0 & 0.001 & 1.21e-10 & 0 & 1.2 \\
100.0 & 0.1   & 5.51e-6  & 0 & 1.1 \\
100.0 & 1.0   & 8.46e-2  & 0 & 2.1 \\
1000.0& 0.001 & 1.40e-9  & 0 & 1.1 \\
1000.0& 0.1   & 7.06e-6  & 0 & 1.1 \\
1000.0& 1.0   & 8.46e-2  & 0 & 2.3 \\
3000.0& 0.001 & 7.44e-12 & 0 & 1.4 \\
3000.0& 0.1   & 7.18e-6  & 0 & 0.9 \\
3000.0& 1.0   & 8.46e-2  & 0 & 2.0 \\
\bottomrule\end{tabular}\end{center}

(Source: \texttt{method\_c\_v2\_grid/grid.json}.) The residual is
essentially $\tau$-only (not $\gamma$); at $\tau = 1$ it floors at
$\|F\|_\infty \approx 0.08$--$0.10$. This residual is the
non-additive correction term --- a structural feature of the FP that
the additive ansatz cannot capture. Aggregated by $\tau$:

\begin{center}\small\begin{tabular}{rrr}\toprule
$\tau$ & max $\|F\|_\infty$ & median $\|F\|_\infty$ \\
\midrule
0.001 & 1.24e-8 & 1.40e-9 \\
0.005 & 1.72e-6 & 3.73e-9 \\
0.01  & 1.37e-5 & 1.71e-8 \\
0.05  & 1.64e-3 & 1.98e-6 \\
0.1   & 1.05e-2 & 1.57e-5 \\
0.3   & 5.91e-2 & 2.25e-3 \\
0.5   & 4.16e-2 & 3.42e-2 \\
0.7   & 6.07e-2 & 6.06e-2 \\
0.85  & 6.67e-2 & 6.65e-2 \\
1.0   & 1.02e-1 & 8.46e-2 \\
\bottomrule\end{tabular}\end{center}

The transition is sharp around $\tau \approx 0.3$: below it the
ansatz is essentially exact (median $\|F\| < 10^{-3}$); above it the
non-additive correction saturates the residual to $\sim 0.05$--$0.10$.

\begin{figure}[H]\centering
  \includegraphics[width=0.95\linewidth]{mc_v2_grid_figs/fig1_F.png}
  \caption{Method C v2 100-cell $(\gamma, \tau)$ grid: $\|F\|_\infty$
    heatmap. The residual depends almost entirely on $\tau$, not
    $\gamma$ (consistent with the additive structure being
    $\tau$-driven).}
\end{figure}

\begin{figure}[H]\centering
  \includegraphics[width=0.95\linewidth]{mc_v2_grid_figs/fig4_cross.png}
  \caption{Cross-sections at fixed $\gamma$ across $\tau$: every
    $\gamma$-row shows the same transition pattern, confirming the
    factorization $\|F\|_\infty(\gamma, \tau) \approx \|F\|_\infty(\tau)$
    in this ansatz.}
\end{figure}

% =====================================================================
\section{The non-additivity at $\tau \to 1$: what the ansatz misses}
% =====================================================================

At $\tau = 0$ the agent signals are uninformative, posteriors are at
prior, and there is no informational coupling. At $\tau = 1$ each
signal carries substantial information; price aggregates them; the
agents' inferences couple through the price (Hellwig's
``noisy rational expectations''). The coupling term is genuinely
non-additive in $(u_1, u_2, u_3)$.

A quantitative way to see this: at $\tau = 1, \gamma = 100$, compute
the residual $R = P_{\rm strict} - \sigma(h(u_1) + h(u_2) + h(u_3))$
under the best-fit Method C v2 $h$. The residual is bounded by
$\sim 0.05$ in magnitude and lives predominantly on the submanifold
where two of the three signals are aligned and the third is opposed
(``the dissenter'' regime). This is exactly where the price aggregator
must work hardest to weigh agreement against dissent.

\textbf{Natural extension:} Method C v3 with a multiplicative or
triple-product correction:
\[
  P \;=\; \sigma\bigl(h(u_1) + h(u_2) + h(u_3) + \kappa \cdot g(u_1, u_2, u_3)\bigr),
\]
where $g$ is an even-degree-2 symmetric polynomial in (centered)
$h(u_i)$ values, e.g. $g = (h(u_1) - h(u_2))(h(u_2) - h(u_3))(h(u_3) - h(u_1))$
or simpler pairwise variances $g = \sum_{i < j} (h(u_i) - h(u_j))^2$.
Estimated 7--9 parameters; expected to drop $\|F\|_\infty$ at $\tau{=}1$
from 0.05 to $< 10^{-3}$ based on the scale of the residual.

% =====================================================================
\section{The geometry of the cusps in the lookup function}
% =====================================================================

At fixed $u_k$ slice, the lookup function $\mu(p, u_k)$ is determined
by the co-area integral over the 1D level set $\{P(u_a, u_b, u_k) = p\}$.
This integral inherits a $|\nabla P|^{-1}$ factor that becomes singular
at critical points of $P$ (i.e. where $\partial_{u_a} P = \partial_{u_b} P = 0$).

\textbf{In the $\tau = 0.001$ regime}: the only critical points are
the corner saddles at $u_a, u_b = \pm U_{\max}$; these are outside the
quadrature domain at any practical $U_{\max}$, so the integral is
bounded and $\mu(p, u_k)$ is $C^\infty$.

\textbf{In the $\tau = 1$ regime}: interior critical points emerge.
On the $G{=}21$ Lobatto slice at $u_k = 0$, we count $\sim 4{-}6$
interior critical points in $(u_a, u_b) \in [-2.33, 2.33]^2$. Each
contributes a logarithmic cusp to $\mu(p, u_k)$ at the corresponding
critical-value $p$. These cusps are the $C^k$ feature that limits
the global Cheby convergence rate.

\textbf{Implication for Option B+.} The empirical CDF, being a
shape-preserving (PCHIP) interpolant rather than a polynomial, handles
the cusps gracefully --- it represents them as kinks in the CDF without
spurious oscillation. This is structural, not numerical: PCHIP is
$C^1$, never $C^\infty$, so it can match the local regularity of
$\mu$ exactly.

\begin{figure}[H]\centering
  \includegraphics[width=0.95\linewidth]{cusps_figs/fig2_grad_critical.png}
  \caption{Critical points of $P(u_1, u_2)$ at $u_3 = 0$ slice,
    $\gamma = 100, \tau = 1$: the zeros of $\nabla P$ live on a
    1D submanifold that crosses the cube along distinct rays. Each
    critical point contributes a logarithmic cusp at the
    corresponding $p$ value.}
\end{figure}

\begin{figure}[H]\centering
  \includegraphics[width=0.95\linewidth]{cusps_figs/fig2_mu_cusps.png}
  \caption{Resulting $\mu(p, u_k)$ cusps: each spike in the
    horizontal-axis derivative corresponds to a critical-point
    crossing. PCHIP-based builders absorb these as kinks; polynomial
    builders ring around them.}
\end{figure}

\begin{figure}[H]\centering
  \includegraphics[width=0.95\linewidth]{cusps_figs/fig3_decay.png}
  \caption{Decay of the Chebyshev coefficients of $P$ on Lobatto
    nodes: $\sim 1/n^{0.8}$ at $\tau = 1$, consistent with $C^0$
    along the ridge submanifold. The slow algebraic decay is the
    quantitative signature of the non-smoothness.}
\end{figure}

% =====================================================================
\section{A digression: the joint density brainstorm}
% =====================================================================

Mid-session we tested a different parametrization: instead of fitting
$\mu(p, u_k)$ directly, fit the joint density $f(p, \mu)$ on
$(p, \mu) \in [0,1]^2$ and recover $\mu(p, u_k) = \mathbb{E}_{f}[\mu \mid p, u_k]$.

The idea: $f$ is smooth (a 2D density on a compact square) while
$\mu(p, u_k)$ inherits the cusps of $|\nabla P|^{-1}$; if we work at
the density level, the cusps become integrable.

\textbf{Why it didn't pan out (in this session).} The joint density
construction would need its own discretization and its own integral
operator; the chain ``cube $\to$ joint density $\to$ marginal $\to$
expectation'' has more moving parts than the direct co-area POU path
and turned out to be no easier to make accurate. We did not invest
deeply because Option B+ was already meeting the median-accuracy bar
with the existing cube discretization.

\textbf{Where this idea might still pay off.} If we ever need
moments of $\mu | p$ beyond the mean (variance, skewness) for some
downstream economic statistic, the joint-density path computes them
all from one fit. Keep in toolbox.

% =====================================================================
\section{A digression: posterior log-odds interpolation}
% =====================================================================

Another mid-session idea: interpolate $L(p, u_k) := \operatorname{logit}\mu(p, u_k) \in
\mathbb{R}$ instead of $\mu \in [0,1]$. Two motivations: (i) the
sigmoid wraps any interpolant back into $[0,1]$ so we never need to
clip, (ii) the linear-Gaussian limit $\tau \to 0$ makes $L$ exactly
linear in $u_k$ and exactly linear in $\operatorname{logit} p$, so
low-degree polynomial fits in $L$-space should track the leading
behavior with very few parameters.

We tested this against the additive ansatz both as a fit target and
as a loss space. The findings, briefly:

\begin{itemize}\itemsep1pt
  \item \textbf{As a fit target,} interpolating $L$ is mildly better
        than interpolating $\mu$ for $|p - 0.5| > 0.3$ but mildly
        worse near $p = 0.5$ (where $L$ is the most sensitive). Net:
        comparable max error, slightly different distribution.
  \item \textbf{As a loss space,} fitting Method C v2 with an
        $\|L_{\rm pred} - L_{\rm strict}\|_\infty$ loss instead of
        an $\|\mu - \mu\|_\infty$ loss did not improve the converged
        $\|F\|_\infty$ (median $1.7 \cdot 10^{-3}$ in both). The
        bottleneck is ansatz expressiveness, not loss geometry.
\end{itemize}

\textbf{Where this idea remains useful.} For storing the lookup table
on disk, $L$-space is more compact (fewer bits per stored value at
target accuracy) and avoids clip-related boundary glitches.

% =====================================================================
\section{Performance comparison at the same $G$}
% =====================================================================

\begin{center}\begin{tabular}{lrrr}\toprule
method & lookup max$|d|$ & median$|d|$ & per-call cost \\
       &                 &              & ($G{=}11$) \\
\midrule
kernel-band R4 (start of session) & 0.5--0.7 & 0.05 & 400 ms \\
Option B+ (cube CDF + PCHIP, numba) & 0.4 & 5e-5 & 50 ms \\
Option B+ refined ($4{\times}$) & 0.2 & 5e-6 & 80 ms \\
Option D (adaptive $p$-grid) & 0.3 & 4e-5 & 55 ms \\
strict-$h{=}0$ (POU + linear, ground truth) & --- & --- & 30 ms \\
Cheby-Lobatto + analytic contour ($G{=}21$) & 0.38 & 0.08 & $\sim 10$ s \\
Method C v2 ansatz, $\tau{=}0.001$ & $1.2 \cdot 10^{-10}$ & $1.4 \cdot 10^{-9}$ & 10 $\mu$s \\
Method C v2 ansatz, $\tau{=}1$ & 0.10 & 0.08 & 10 $\mu$s \\
\bottomrule\end{tabular}\end{center}

% =====================================================================
\section{Honest assessment and where to invest next}
% =====================================================================

\subsection{What we actually achieved}

\begin{enumerate}\itemsep1pt
  \item We diagnosed and quantified the kernel-band-Richardson bias.
        The earlier ``deficit ridge'' result needs revision: under the
        correct operator, the deficit collapses $\sim 40{\times}$.
  \item We built two production-quality lookup evaluators (Option B+
        and Option D) that achieve median accuracy $\sim 10^{-5}$
        --- four decades better than what we had at the start of the
        session --- at competitive cost.
  \item We characterized the smoothness barrier (algebraic Cheby
        convergence) and identified its geometric source (interior
        critical points of $P$ at high $\tau$).
  \item We have a closed-form-exact ansatz at $\tau \to 0$ that
        solves a 100-cell $(\gamma, \tau)$ grid in 2 min and a
        1000-cell $\tau$-ladder in 17 min with zero monotonicity
        violations.
  \item We have an analytic 1D-integral lookup ($\sim 10\,\mu$s per call)
        for the additive-ansatz regime.
\end{enumerate}

\subsection{What we did not achieve}

\begin{enumerate}\itemsep1pt
  \item Machine-$\varepsilon$ lookups in the sharp ($\tau \to 1$) regime.
        Cheby spectral can't get there without piecewise patches; B+
        floors at $\sim 10^{-5}$ on the median and $\sim 10^{-1}$ on
        the max.
  \item A FP solver that combines the structural ansatz (cheap, smooth,
        symmetry-clean) with the non-additive correction (essential at
        $\tau$ near 1). Method C v3 is the obvious next step but not
        implemented.
  \item A Cheby-Hermite multi-element spectral solver. Big code, would
        give exponential convergence per patch but adds patch-matching
        complications.
\end{enumerate}

\subsection{Recommendations in priority order}

\begin{enumerate}\itemsep1pt
  \item \textbf{Switch production lookup to Option B+ numba} for all
        $(\gamma, \tau)$ cells. Drop-in replacement; $\sim 100{\times}$
        median-accuracy gain over kernel-band R4 at lower wall.
  \item \textbf{Recompute the deficit ridge plot} from the K=3 paper
        using strict-$h{=}0$ FPs. The qualitative pattern likely
        survives (we showed the deficit is still non-monotone in
        $\tau$), but quantitative values shift by $40{\times}$ in
        most cells.
  \item \textbf{Implement Method C v3} (additive + non-additive
        correction). Expected payoff: a closed-form-near FP solver
        across the full $(\gamma, \tau)$ grid in $< 1\,\mathrm{min}$,
        which would let us run continuation studies (e.g. across
        the fold at $\tau \approx 0.79$) in interactive time.
  \item \textbf{Multi-element Chebyshev} only if pursuit of
        machine-$\varepsilon$ machine-eps lookups becomes a binding need
        --- this is a research project, not a quick fix.
  \item \textbf{Skip} further investment in kernel-band Richardson:
        bias is fundamental.
\end{enumerate}

% =====================================================================
\section{Algorithm sketches}
% =====================================================================

\subsection{Strict-$h{=}0$ co-area + POU lookup}

\begin{verbatim}
def mu_strict(P_grid, u_grid, p_target, u_k, tau):
    # 1. find all cube cells whose [P_min, P_max] straddles p_target
    cells = find_straddle_cells(P_grid, p_target, u_k)
    A0, A1 = 0.0, 0.0
    for cell in cells:
        # 2. parametrize the level-set arc through the cell via
        #    bilinear inversion (closed-form roots of bilinear)
        arc = bilinear_level_arc(P_grid[cell], p_target)
        # 3. apply POU weight along the arc (Cinfty bump)
        pou = pou_weight(arc, cell)
        # 4. integrate f_v(u_a)*f_v(u_b) along the arc
        a0, a1 = co_area_integral(arc, pou, tau)
        A0 += a0; A1 += a1
    f0k = gauss(u_k, -0.5, tau); f1k = gauss(u_k, +0.5, tau)
    return (f1k * A1) / (f0k * A0 + f1k * A1)
\end{verbatim}

\subsection{Option B+ cube CDF + PCHIP density lookup}

\begin{verbatim}
def mu_optB(P_grid, u_grid, p_grid, u_k, tau, refine=4):
    # 1. for each cube cell, compute Gaussian mass under (f0, f1)
    M0, M1 = gaussian_cell_mass(u_grid, tau, refine)
    # 2. build empirical CDF in p: F0(p), F1(p) = sum of cell mass
    #    over cells with P_cell <= p, weighted by tau-Gaussian
    F0 = empirical_cdf(M0, P_grid, p_grid, u_k)
    F1 = empirical_cdf(M1, P_grid, p_grid, u_k)
    # 3. PCHIP-differentiate to get density f0(p|u_k), f1(p|u_k)
    f0p = pchip_derivative(p_grid, F0)
    f1p = pchip_derivative(p_grid, F1)
    # 4. ratio
    return f1p / (f0p + f1p)
\end{verbatim}

\subsection{Method C v2 ansatz analytic 1D-integral lookup}

\begin{verbatim}
def mu_methodC_v2(coefs, p, u_k, tau, NQK=24):
    # h(u) = integral of sigma^2(s) ds from 0 to u, sigma^2 even Cheby
    h_uk = h_at(u_k, coefs)
    L = logit(p)
    # Level set in (h(u_a), h(u_b)) is the affine line v_a + v_b = L - h(u_k)
    v_lo = max(h_min, L - h_uk - h_max)
    v_hi = min(h_max, L - h_uk - h_min)
    # 1D GL quadrature in v_a
    xi, w = gauss_legendre(NQK)
    v_a = 0.5*(v_lo+v_hi) + 0.5*(v_hi-v_lo)*xi
    v_b = L - h_uk - v_a
    # invert h to get u_a, u_b
    u_a = h_inverse(v_a, coefs); u_b = h_inverse(v_b, coefs)
    # integrand: gaussian densities / (sigma^2 jacobian)^2
    sa = sigma2_at(u_a, coefs); sb = sigma2_at(u_b, coefs)
    f0a = gauss(u_a, -0.5, tau); f1a = gauss(u_a, +0.5, tau)
    f0b = gauss(u_b, -0.5, tau); f1b = gauss(u_b, +0.5, tau)
    A0 = sum(w * f0a * f0b / (sa * sb))
    A1 = sum(w * f1a * f1b / (sa * sb))
    f0k = gauss(u_k, -0.5, tau); f1k = gauss(u_k, +0.5, tau)
    return (f1k * A1) / (f0k * A0 + f1k * A1)
\end{verbatim}

The structural difference between the three: strict-$h{=}0$ traces the
level set explicitly; B+ accumulates cube mass and differentiates
empirically; Method C v2 inverts the ansatz's $h$ to give a closed-form
1D integral. They span a spectrum of how much structure the operator
exploits.

% =====================================================================
\section{The fold at $\tau \approx 0.79$}
% =====================================================================

A practical bound on the descent: pseudo-arclength continuation in
$\tau$ from $\tau = 2.0$ downward (DD solver, all 32 digits) shows the
minimum singular value of $I - J$ at the FP collapsing:
\[
  \sigma_{\min}(I - J) :\;\; 0.45\; (\tau=1.0)\;\; \to\;\; 0.31\; (\tau=0.9)\;\; \to\;\; 0.017\; (\tau=0.8),
\]
suggestive of a saddle-node fold near $\tau \approx 0.79$. Below the
fold the operator likely admits multiple FPs (or a continuous family);
the Newton-based descent we use cannot bridge the fold. This is a
fundamental feature of the equilibrium, not a numerical glitch.

\textbf{Implication for the $\tau$-ladder.} The 1000-cell ladder
($\tau \in [10^{-3}, 1.0]$) is on the smooth branch above the fold.
Method C v2's TRF fit converges across that entire branch with zero
monotonicity violations. Extending below $\tau = 0.79$ would require
pseudo-arclength continuation (not branch chasing in $\tau$ alone)
and is a substantial separate exercise.

\textbf{Below-fold candidates not pursued in this session:}
\begin{itemize}\itemsep1pt
  \item Pseudo-arclength with a path-following predictor.
  \item Branch enumeration via homotopy from the trivial $P \equiv 0.5$
        FP (which is always a FP by sign-flip symmetry).
  \item Continuation in $\gamma$ across $\tau$ slices below the fold.
\end{itemize}

% =====================================================================
\section{Open question: how non-additive is the FP exactly?}
% =====================================================================

The Method C v2 residual at $\tau = 1, \gamma = 100$ is bounded by
$\sim 0.05$ in max-norm and concentrates on the dissenter regime
(two signals aligned, one opposed). Three structural conjectures:

\begin{enumerate}\itemsep1pt
  \item \textbf{Conjecture A:} the non-additivity is exactly cubic in
        the centered signals,
        $g(u_1,u_2,u_3) = \kappa(\tau) \cdot (h(u_1)\, h(u_2)\, h(u_3))$
        with $\kappa(0) = 0$ and $\kappa(\tau) = O(\tau^2)$. Test:
        $\kappa$ should fit cleanly to a power law in $\tau$; if it
        does, the additive + triple-product ansatz nails the FP
        across the entire above-fold branch.
  \item \textbf{Conjecture B:} the non-additivity is pairwise
        variance,
        $g(u_1,u_2,u_3) = \kappa(\tau) \cdot \sum_{i<j} (h(u_i) - h(u_j))^2$.
        Same parameter count; different functional form. Distinguished
        from Conjecture A by the sign of the residual at the
        ``all-equal'' configuration $u_1 = u_2 = u_3$ (A says zero, B
        says positive).
  \item \textbf{Conjecture C:} the non-additivity contains a kink
        --- not a smooth correction, but a piecewise-defined term that
        switches on at a $\tau$-dependent threshold. This would
        explain why the Method C v2 residual is concentrated rather
        than smoothly distributed, and would predict that no smooth
        finite-parameter correction can drive $\|F\|_\infty$ below a
        $\tau$-dependent floor.
\end{enumerate}

We have not run the experiment that distinguishes among A, B, C.
Quickest test: at $\tau = 1, \gamma = 100$, evaluate the strict-$h{=}0$
FP residual on the diagonal $u_1 = u_2 = u_3 = s$ and off-diagonal
slices; the functional form will fingerprint the conjecture in
$\sim 100$ FP evaluations.

% =====================================================================
\section{Convergence rates: a stylized derivation}
% =====================================================================

To understand \emph{why} the various operators show the convergence
rates they do, it helps to write down the leading error term in each
case.

\subsection{Kernel-band Richardson on a $C^k$ integrand}

A standard Richardson argument assumes the kernel-band lookup admits
an asymptotic expansion in $h^2$:
\[
  \mu_h(p, u_k) \;=\; \mu_0(p, u_k) + c_1 h^2 + c_2 h^4 + \cdots
\]
With $R$ band evaluations at $h_1, \ldots, h_R$, an order-$R$
Vandermonde extrapolation kills the first $R$ terms and leaves
$O(h_{\max}^{2(R+1)})$. \emph{This requires the integrand to be
analytic in $h^2$.}

In our problem, the level set $\{P = p\}$ has kinks at critical
points of $P$. The kernel band picks up or drops a kink
discontinuously as $h$ crosses certain values; the contribution of
that kink to the integral is $O(h^\alpha)$ with $\alpha$ \emph{not}
necessarily an even integer (typically $\alpha = 1$ or $3/2$
depending on the local geometry). Richardson in even powers of $h$
cannot annihilate these terms. Empirically the residual floors at the
$O(h^\alpha)$ level rather than at $O(h^{2(R+1)})$, which is exactly
what the smaller-$h$ test (Finding 1) shows.

\subsection{Cube CDF + PCHIP density}

The cube empirical CDF $F_h(p)$ has error $O(G^{-2})$ in $p$ from the
piecewise-bilinear treatment of cube cells (since the cube is a
$G^3$-vertex grid, mass within each cell is computed by Gaussian
integral against a bilinear $P$). Differentiating via PCHIP (which is
$C^1$ and shape-preserving) inherits the $O(G^{-2})$ rate without
introducing additional error in smooth regions. Where $\mu(p, u_k)$
has a kink, PCHIP represents it as a kink in the density (no
oscillation), giving locally $O(\Delta p)$ accuracy --- still
acceptable for the lookup integral.

Predicted rate: $O(G^{-2})$ globally, with $O(\Delta p)$ around cusps.
The Option B+ refinement (sub-grid factor $4{\times}$) is effectively
$G_{\rm eff} = 4G$, hence $\sim 16{\times}$ accuracy gain. Observed:
$\sim 10{\times}$ accuracy gain at refine $= 4$, consistent with the
prediction modulo the cusp regions.

\subsection{Chebyshev on a $C^k$ function}

Chebyshev expansion of a function with $k$ continuous derivatives but
not $C^{k+1}$ has coefficients decaying as $|a_n| \sim n^{-(k+1)}$
(folklore: this is the ``Jackson--Bernstein'' rate). The interpolant
error is then $O(G^{-k})$, algebraic in $G$.

If the function is analytic in a Bernstein ellipse of radius $\rho > 1$
around $[-1, 1]$, coefficients decay as $|a_n| \sim \rho^{-n}$ and
the error is $O(\rho^{-G})$ --- exponential.

Our observed edge-coefficient decay $\sim 1/G^{0.8}$ suggests $k \approx 0$
in the global Chebyshev sense, i.e. the FP is barely $C^0$ along
some ridge submanifold. Investing in higher $G$ on a global Cheby fit
buys very little.

\subsection{Method C v2 ansatz}

The ansatz is a finite-parameter (5 numbers) approximation in
$L^2(\text{cube})$. Its error is bounded by the projection error of
the true FP onto the span of $\{\text{additive sigmoids}\}$, which is
$\tau$-dependent: zero at $\tau = 0$ (the linear-Gaussian limit is
exactly additive), $O(\tau^2)$ at small $\tau$, and $O(0.05)$ saturating
near $\tau = 1$.

The TRF solver itself converges \emph{exponentially} in iterations
(quadratic basin of attraction) and hits the ansatz floor in
$\sim 10$--$30$ iterations.

% =====================================================================
\section{Appendix A: timing breakdown per operator}
% =====================================================================

All numbers below at $G = 11$ DD on a 6-core machine, single-call (not
warmed):

\begin{center}\begin{tabular}{lrrr}\toprule
operator stage & wall (ms) & \% of total & notes \\
\midrule
\multicolumn{4}{l}{\textbf{kernel-band Lin-CDF R4 lookup builder ($\sim$400 ms total)}} \\
~~kernel-band $h=0.5$ pass         & 95  & 24\% & $\mathrm{NQK}{=}16$, $G_p{=}121$ \\
~~kernel-band $h=0.4$ pass         & 95  & 24\% & \\
~~kernel-band $h=0.3$ pass         & 95  & 24\% & \\
~~kernel-band $h=0.2$ pass         & 95  & 24\% & \\
~~Vandermonde solve + assembly     & 20  & 5\%  & 4-point Richardson \\
\midrule
\multicolumn{4}{l}{\textbf{strict-$h{=}0$ co-area POU lookup builder ($\sim$30 ms total)}} \\
~~cube-cell co-area integral       & 18  & 60\% & per $(p, u_k)$ bin \\
~~POU weight assembly              & 7   & 23\% & 3-chart partition \\
~~$\mu$ normalize + clip           & 5   & 17\% & \\
\midrule
\multicolumn{4}{l}{\textbf{Option B+ numba lookup builder ($\sim$50 ms total)}} \\
~~cube-CDF accumulation (jit)      & 25  & 50\% & sum cube-cell mass per $p$-bin \\
~~PCHIP density fit                & 18  & 36\% & shape-preserving derivative \\
~~$\mu$ ratio + normalize          & 7   & 14\% & \\
\midrule
\multicolumn{4}{l}{\textbf{Method C v2 ansatz lookup ($\sim$10 $\mu$s per call)}} \\
~~$h^{-1}$ table lookup            & 3   & 30\% & 201-point cached \\
~~1D GL integral                   & 5   & 50\% & 24 nodes \\
~~ratio + clip                     & 2   & 20\% & \\
\bottomrule\end{tabular}\end{center}

% =====================================================================
\section{Appendix B: per-cell sample lookup values}
% =====================================================================

For a transparent cross-check, here are $\mu(p, u_k = 0)$ at
$\gamma = 100$, $\tau = 1.0$ for a handful of $p$ values, computed by
each major operator family:

\begin{center}\small\begin{tabular}{rrrrr}\toprule
$p$ & R4 (start) & strict $h{=}0$ & Option B+ & Method C v2 \\
\midrule
0.05 & 0.071 & 0.043 & 0.044 & 0.045 \\
0.10 & 0.118 & 0.085 & 0.085 & 0.086 \\
0.20 & 0.207 & 0.166 & 0.166 & 0.168 \\
0.30 & 0.302 & 0.249 & 0.249 & 0.252 \\
0.40 & 0.398 & 0.347 & 0.346 & 0.349 \\
0.50 & 0.500 & 0.500 & 0.500 & 0.500 \\
0.60 & 0.602 & 0.653 & 0.654 & 0.651 \\
0.70 & 0.698 & 0.751 & 0.751 & 0.748 \\
0.80 & 0.793 & 0.834 & 0.834 & 0.832 \\
0.90 & 0.882 & 0.915 & 0.915 & 0.914 \\
0.95 & 0.929 & 0.957 & 0.956 & 0.955 \\
\bottomrule\end{tabular}\end{center}

The R4 column is systematically pulled toward $p$ (the ``identity-like''
bias of the kernel-band operator); strict, B+, and Method C v2 agree
to $\sim 10^{-3}$ across the $p$ grid.

% =====================================================================
\section{Appendix C: monotonicity diagnostic}
% =====================================================================

Among the R4 fixed points used in earlier work, 27\% have at least one
cube cell where $P(u_1, u_2, u_3)$ violates axis-monotonicity in some
$u_i$. This is an operator artifact: the R4 operator does not enforce
monotonicity and the violations live in cells where the kernel band
straddles a critical point of the implicit level set.

Under strict-$h{=}0$ + POU, the violation rate drops to 4\% --- not
zero (the linear interpolation between cube nodes can still violate
monotonicity if adjacent node values cross), but qualitatively
different. Under Method C v2's structural ansatz, the violation
rate is exactly 0\% by construction.

This is a quality dimension distinct from $\|F\|_\infty$: a small
operator residual at a non-monotone FP still represents an invalid
economic equilibrium. We did not enforce monotonicity in any
operator-level solver; for production it is worth adding a
monotonicity-projection step.

% =====================================================================
\section{Appendix D: experiments by date (cumulative log)}
% =====================================================================

This summary distilled $\sim$30 discrete experiments. For
traceability, the rough chronological order:

\begin{enumerate}\itemsep1pt
  \item Auto-NQ-finder + N=8 NQ=14 phase diagram; phase boundary at slope $0.384$.
  \item DD K=3 sweep across $(\gamma, \tau)$: lookup-table architecture.
  \item Variant comparator (V1--V7): max cross-variant spread $\sim 10^{-5}$.
  \item Richardson U-curve R2--R7: R4 minimizes self-residual by construction.
  \item Smaller-$h$ test: bias persists at $h = 0.02$.
  \item Strict-$h{=}0$ cross-validation: $\sim 0.05$--$0.6$ R4 bias confirmed.
  \item Strict-$h{=}0$ FP solve from R4 warm-start: deficit collapses $40{\times}$.
  \item Alternative builders A--E sweep: Option B+ wins on median.
  \item Option B+ numba refinement + BC extrapolation: median $\to 10^{-5}$.
  \item Option B+ as inner FP operator: fails (lacks Lipschitz contraction).
  \item Option D adaptive $p$-grid (4 monitors): arclength wins, $-25\%$ max error.
  \item Cheby $3$D on CDF-uniform: Runge phenomenon, divergent at degree $\ge 13$.
  \item Cheby-Lobatto FP solve at $G = 11, 15, 21$: $\|F\| \sim 10^{-14}$ each.
  \item Cheby-Lobatto lookup convergence: algebraic only ($\sim G^{-1.6}$).
  \item Cheby-native FP (no kernel, no Richardson): confirms algebraic rate is intrinsic.
  \item Joint $(p, \mu)$ density brainstorm: did not pan out, kept in toolbox.
  \item Logit-space loss for ansatz: no improvement vs $\mu$-space.
  \item Method C v1 (7-param): could not separate signal-strength from coupling.
  \item Method C v2 (5-param, $\sigma^2$ Cheby): near-exact at $\tau = 0.001$ ($\|F\| \sim 10^{-10}$ across $\gamma$).
  \item Analytic 1D-integral lookup under Method C v2: $\sim 10\,\mu$s per call.
  \item 100-cell $(\gamma, \tau)$ grid via Method C v2: 2 min, zero violations.
  \item 1000-cell $\tau$-ladder via Method C v2: 17 min, zero violations.
  \item Method C v2 with strict-$h{=}0$ as inner operator: trivial FP, multiple solutions.
  \item Monotonicity diagnostic: 27\% violation rate under R4, 4\% under strict $h{=}0$.
  \item Fold detection: $\sigma_{\min}(I-J) = 0.017$ at $\tau = 0.8$.
\end{enumerate}

Not all experiments produced positive results; the negative ones (B+
as solver, Cheby on CDF-uniform, joint density, Method C v1) are
documented here so we don't repeat them.

% =====================================================================
\section{Appendix E: a one-page through-line}
% =====================================================================

\textbf{Where we started.} A kernel-band Lin-CDF R4 Richardson lookup
on a CDF-uniform $G{=}11$ DD-precision $u$-grid. Self-residual
$|F| \sim 10^{-14}$ at its own FP; we thought this meant the operator
was correct. It did not.

\textbf{What we discovered.} The R4 lookup is systematically biased
$\sim 0.05$--$0.6$ in $\mu$ versus the strict-$h{=}0$ co-area + POU
reference, depending on regime. Richardson cannot extrapolate to
$h = 0$ because the level-set integral is not analytic in $h^2$ ---
kinks at critical points of $P$ inject $h^\alpha$ terms with
non-integer $\alpha$.

\textbf{What it cost us economically.} At $\gamma = 100, \tau = 1$
the R4 deficit was $0.255$. Under strict $h = 0$ the deficit is
$0.0065$ --- a $40{\times}$ collapse. The ``deficit ridge'' that
motivated earlier work was largely a kernel-band-Richardson artifact.

\textbf{What we built to replace it.}
\begin{itemize}\itemsep1pt
  \item Option B+ (cube empirical CDF + PCHIP density): median lookup
        error $\sim 10^{-5}$ at $\sim 50$ ms per call. Drop-in
        replacement for the kernel-band lookup. Recommended for
        production.
  \item Option D (adaptive $p$-grid by arclength): stacks with B+ for
        $\sim 25\%$ further max-error reduction at no cost.
  \item Method C v2 (structural ansatz $P = \sigma(\sum_i h(u_i))$):
        $\sim 10\,\mu$s per call. Exact in the $\tau \to 0$ limit;
        $\|F\|_\infty \sim 0.05$ at $\tau = 1$ (non-additive correction
        needed; Method C v3 is the next step).
\end{itemize}

\textbf{What we ruled out.}
\begin{itemize}\itemsep1pt
  \item Richardson on smaller $h$: bias persists at $h = 0.02$.
  \item Cheby on CDF-uniform: Runge phenomenon, divergent.
  \item Cheby-Lobatto + analytic contour: algebraic-only convergence,
        not competitive (the FP is not analytic).
  \item Joint $(p, \mu)$ density: more moving parts, no easier than
        direct co-area.
  \item Logit-space loss for Method C v2: no improvement.
  \item Method C v1 (7-parameter): could not separate signal-strength
        from agent-coupling; rejected for v2's 5-parameter
        $\sigma^2(s)$ formulation.
\end{itemize}

\textbf{What remains open.}
\begin{itemize}\itemsep1pt
  \item Machine-$\varepsilon$ lookups at high $\tau$: needs multi-element
        spectral or kink-aware quadrature.
  \item Below-fold ($\tau < 0.79$) continuation: needs pseudo-arclength.
  \item Method C v3 (non-additive correction term): expected to nail
        the high-$\tau$ regime; not implemented this session.
  \item Whether the K=3 paper's deficit-ridge conclusions need
        retraction or just quantitative revision.
\end{itemize}

\textbf{The single most important conclusion of the session.} The
self-residual of an operator on its own FP is \emph{not} a measure of
correctness. A biased operator can have arbitrarily small self-residual
because the FP migrates to compensate for the bias. The right metric is
agreement with an independently-derived ground-truth operator (here:
strict $h = 0$). This is a generic lesson, not specific to K=3.

% =====================================================================
\section{Reproducibility}
% =====================================================================

\subsection{Source files (in \texttt{projects/REZN/solver\_code/highprec\_dd\_qd/})}

\begin{itemize}\itemsep1pt
  \item \texttt{dd\_ops.py}: double-double arithmetic primitives.
  \item \texttt{dd\_solver.py}: outer Anderson FP solver with PCHIP $\mu$.
  \item \texttt{dd\_k3\_solver.py}: K=3 specific operator.
  \item \texttt{dd\_k3\_strict\_solve.py}: strict-$h{=}0$ co-area + POU solver.
  \item \texttt{dd\_k3\_compare.py}: 7-variant lookup comparator.
  \item \texttt{dd\_k3\_rconverge.py}: Richardson order U-curve.
  \item \texttt{dd\_k3\_small\_h.py}: smaller-$h$ test.
  \item \texttt{dd\_k3\_lookup\_alts.py}: alternative builders A--E.
  \item \texttt{dd\_k3\_optB\_numba.py}, \texttt{dd\_k3\_optB\_refined.py}: Option B+ production.
  \item \texttt{dd\_k3\_optD\_adaptive\_p.py}: Option D adaptive $p$-grid.
  \item \texttt{dd\_k3\_cheby\_h0.py}, \texttt{dd\_k3\_cheby\_lobatto.py}:
        Chebyshev brainstorm (CDF-uniform failure + Lobatto convergence).
\end{itemize}

Method C v2 source (was being developed in \texttt{/tmp} and lost to
the container restart; reconstructible from this report, $\sim 200$ LOC).

\subsection{Data}

\texttt{projects/REZN/solved\_fixed\_points/dd\_k3\_overnight/} contains
all JSON results, NPZ fixed points, and PNG figures used in this report.

\subsection{Standalone reports also in repo}

\begin{itemize}\itemsep1pt
  \item \texttt{dd\_k3\_overnight\_report.pdf}: kernel-band Richardson
        diagnostic (24p).
  \item \texttt{dd\_k3\_alts\_report.pdf}: alternative builders A--E (19p).
  \item \texttt{optB/dd\_k3\_optB\_final\_report.pdf}: Option B+
        production write-up (20p).
  \item \texttt{optD/dd\_k3\_optD\_report.pdf}: Option D adaptive $p$-grid
        (12p).
  \item \texttt{cheby/dd\_k3\_cheby\_report.pdf}: Chebyshev + analytic
        contour brainstorm (4p).
\end{itemize}

This summary is the consolidated through-line; the standalone reports
contain additional detail and the underlying ground-truth tables.

% =====================================================================
\section{Sanity-check figures: lookups across operators}
% =====================================================================

The figures below were generated during the alternative-builder sweep
and re-used here to visually anchor the quantitative claims about
operator agreement.

\begin{figure}[H]\centering
  \includegraphics[width=0.95\linewidth]{alts_figs/fig2_mu_curves.png}
  \caption{Lookup $\mu(p, u_k)$ across builders, one $u_k$ slice per
    panel. Strict-$h{=}0$ (black) vs Option B+ (blue) overlap to
    plotting accuracy; kernel-band R4 (red) is visibly offset; Option
    A (Cheby on CDF-uniform) shows tail Runge oscillations.}
\end{figure}

\begin{figure}[H]\centering
  \includegraphics[width=0.95\linewidth]{alts_figs/fig4_distribution.png}
  \caption{Histogram of $|\mu_{\rm method} - \mu_{\rm strict}|$ across
    all $(p, u_k)$ pairs at $G = 11$, $(\gamma, \tau) = (100, 1)$. The
    Option B+ distribution sits two decades to the left of all the
    kernel-band variants.}
\end{figure}

% =====================================================================
\section{Acknowledgements and a note on uncertainty}
% =====================================================================

Every quantitative claim in this report is rooted in a specific
operator at a specific discretization. Where two operators disagree
(e.g. R4 vs strict $h{=}0$), we have generally trusted the operator
whose construction relies on fewer extrapolation assumptions (strict
$h{=}0$ over Richardson), even though this means accepting larger
$\|F\|_\infty$ on the trusted side. This is a methodological choice
rather than a proof.

Three places where additional care would be warranted before any of
this is used downstream:

\begin{enumerate}\itemsep1pt
  \item \textbf{The deficit ridge.} Our conclusion that the ridge is
        largely an artifact rests on strict-$h{=}0$ FPs with
        $\|F\|_\infty \sim 10^{-3}$, not $\sim 10^{-14}$. The
        \emph{qualitative} reduction by $40{\times}$ is robust to
        order-of-magnitude changes in the strict-$h{=}0$ residual,
        but the \emph{exact} new deficit values could shift by a
        factor of $\sim 2$ as the strict solver is refined.
  \item \textbf{The fold location.} $\tau \approx 0.79$ is the value
        at which Newton stops finding the same FP under the operators
        and warm-starts we tried. It is consistent with a true fold
        but a continuation study would be required to verify.
  \item \textbf{The non-additive correction.} We have evidence that
        the additive ansatz misses $\sim 0.05$ in $\|F\|_\infty$ at
        $\tau = 1$, but we have not validated whether this is a
        smooth correction (Conjectures A, B) or a kink-type
        non-smoothness (Conjecture C). The Method C v3 ansatz design
        depends on which it is.
\end{enumerate}

\end{document}
